\documentclass{situationist-preprint}

\title{The Effectiveness of Mathematics in Science}
\msid{SCI-17-a}
\msfamily{SCI-17}
\msdiscipline{Philosophy of Science}
\mstype{Principal manuscript}
\mskeywords{applicability of mathematics; Wigner; unreasonable effectiveness; selection; Steiner; structuralism}

\begin{document}
\maketitle

\begin{abstract}
\citet{Wigner1960} called the effectiveness of mathematics in the natural sciences unreasonable: concepts developed for their beauty turn out, much later, to describe the world with astonishing accuracy. We argue that the effectiveness is not unreasonable, because scientists select, from a very large stock of mathematical structures, those that fit the phenomena. The surprise Wigner reports depends on counting the successes of applied mathematics without the much larger number of structures never applied and phenomena never mathematized. Much of the mathematics later applied was, moreover, abstracted from structures of geometric or physical origin. We answer the objections from unexpected precision, from Steiner's formal analogies and from dispensability. Selection explains why the mathematics in use fits, given that mathematical structures can be instantiated by physical systems; it does not explain applicability away, and it is compatible with the indispensability of mathematics to physics.
\end{abstract}

\section{The problem}

Wigner's examples are well known. Complex numbers, introduced to solve cubic equations, are the language of quantum mechanics; matrix algebra, developed for its own sake, gave Heisenberg his mechanics; the law of gravitation, stated on the basis of rough observation, proved accurate to far better than the data that suggested it \citep{Wigner1960}. \citet{Hamming1980} added that mathematics is effective even where we would not expect it, and offered partial explanations, among them that we see what we look for and select the mathematics that fits, while finding them insufficient. \citet{Steiner1998} argued that the success of mathematical analogies in physical discovery, as when physicists predicted particles by extending the symmetry groups that classified known ones, shows that the universe is in a sense \enquote{user-friendly} to human mathematical aesthetics, a conclusion he took to tell against naturalism.

\section{The argument}

\begin{definition}[Effectiveness]\label{d:eff}
Mathematics is \emph{effective} in a science when some of the structures it studies are instantiated, to good approximation, by the systems the science describes, so that calculation within those structures yields reliable predictions about the systems.
\end{definition}

\begin{definition}[Unreasonable]\label{d:unr}
The effectiveness of mathematics is \emph{unreasonable} if the proportion of mathematical structures that turn out effective is much higher than we have reason to expect, given how they were developed.
\end{definition}

\begin{assumption}[Abundance]\label{a:abund}
Mathematics studies a very large range of structures, most of which have never been applied in any science.
\end{assumption}

\begin{assumption}[Selection]\label{a:sel}
Scientists search the available structures for those that fit a given range of phenomena, retain those that fit and discard those that do not; and phenomena that no available structure fits are not counted among the domains to which mathematics applies.
\end{assumption}

\begin{theorem}\label{t:main}
The effectiveness of mathematics is not unreasonable.
\end{theorem}

\begin{proof}
Wigner's examples are cases in which a structure turned out to fit. The relevant proportion, by Definition~\ref{d:unr}, is that of effective structures among all structures developed. By Assumption~\ref{a:abund}, the denominator is very large, and the effective structures are a small fraction of it: most of group theory, topology, number theory and analysis has found no physical application. By Assumption~\ref{a:sel}, the effective structures we observe are those that survived a search for fit, so that their fit is the criterion of their selection, not a coincidence among unselected samples; and the phenomena we observe to be mathematically tractable are those for which a fitting structure was found, while turbulence, much of biology and most of social life resist precise mathematical treatment. The observed successes are therefore what a large stock of structures and a selective search would produce, and there is no proportion of successes that exceeds what we have reason to expect.
\end{proof}

The argument needs the premise that mathematical structures can be instantiated by physical systems, which is presupposed by Definition~\ref{d:eff}. It does not explain applicability from nothing, and it does not require that mathematics be a game of symbols. On the contrary, MTH-15-a argues that formalism cannot explain applicability, since applicability requires that calculation preserve a structure shared by the mathematics and its subject matter; and the structuralism of MTH-14, on which numbers are positions in structures, gives the selection argument the stock it draws on. Selection explains which of the many structures are applied, and why those applied fit, given that structures can be instantiated.

A second factor reduces the surprise further. Much of the mathematics later applied was abstracted from structures that already had physical or geometric origins. Complex numbers generalize rotations and scalings in the plane; group theory generalizes the symmetries of figures and equations; Riemannian geometry generalizes the intrinsic geometry of surfaces. That generalizations of structures abstracted from the world should again describe aspects of the world is not surprising \citep{GrattanGuinness2008}.

\section{Objections}

\begin{objection}[Accuracy]
Selection explains why mathematics fits the data used to select it, not why it fits new data with greater precision than the original evidence, as Newton's law did.
\end{objection}

\begin{reply}
The extra precision is evidence that the selected structure is approximately instantiated by the system, not merely fitted to a sample of its behavior. That a structure selected on rough data turns out to be instantiated is the ordinary success of theory, which SCI-02 explains by the approximate truth of the theory's essential posits; it requires no further fact about mathematics.
\end{reply}

\begin{objection}[Steiner's analogies]
In the cases \citet{Steiner1998} describes, physicists did not search for a structure that fit; they extended a structure for reasons of mathematical form, and nature complied.
\end{objection}

\begin{reply}
Heuristics based on mathematical form are themselves selected: physicists use the analogies, such as symmetry principles, that have worked, and the failures of formal analogy are rarely recorded \citep{Colyvan2001}. Symmetry principles, moreover, are not merely aesthetic: they encode conservation laws and invariances that experience had already shown to hold widely. Selection operates on strategies as well as on structures.
\end{reply}

\begin{objection}[Dispensability]
If mathematics is effective only because it is selected for fit, it is a convenient tool, and science could in principle do without it.
\end{objection}

\begin{reply}
Selection explains why the mathematics we use fits; it says nothing about whether the same content could be stated without it. MTH-16-a argues that it cannot, since physics cannot be nominalized, and the selection account is compatible with that conclusion: what is selected is a structure that the physical system instantiates, and the theory's claims are partly about that structure.
\end{reply}

\section{Conclusion}

The effectiveness of mathematics in science is not unreasonable. Mathematics offers a vast stock of structures, scientists select those that fit the phenomena, and the domains that resist mathematization are not counted. The successes Wigner found miraculous are what selection from abundance would produce.

\begin{thebibliography}{9}\small
\bibitem[Colyvan(2001)]{Colyvan2001} Colyvan, M. (2001). The miracle of applied mathematics. \emph{Synthese}, 127(3), 265--277.
\bibitem[Grattan-Guinness(2008)]{GrattanGuinness2008} Grattan-Guinness, I. (2008). Solving Wigner's mystery: The reasonable (though perhaps limited) effectiveness of mathematics in the natural sciences. \emph{The Mathematical Intelligencer}, 30(3), 7--17.
\bibitem[Hamming(1980)]{Hamming1980} Hamming, R.~W. (1980). The unreasonable effectiveness of mathematics. \emph{The American Mathematical Monthly}, 87(2), 81--90.
\bibitem[Steiner(1998)]{Steiner1998} Steiner, M. (1998). \emph{The Applicability of Mathematics as a Philosophical Problem}. Harvard University Press.
\bibitem[Wigner(1960)]{Wigner1960} Wigner, E.~P. (1960). The unreasonable effectiveness of mathematics in the natural sciences. \emph{Communications on Pure and Applied Mathematics}, 13(1), 1--14.
\end{thebibliography}

\end{document}
